% Build with: tectonic paper.tex
\documentclass[11pt,letterpaper]{article}
\usepackage[T1]{fontenc}
\usepackage{lmodern}
\usepackage[margin=1in,headheight=15pt]{geometry}
\usepackage{amsmath,amssymb,amsthm}
\usepackage{microtype,booktabs,enumitem,xcolor,fancyhdr,hyperref}
\usepackage[most]{tcolorbox}
\usepackage{array}
\usepackage{needspace}
\newcolumntype{P}[1]{>{\raggedright\arraybackslash}p{#1}}
\definecolor{ink}{HTML}{203B55}
\definecolor{accent}{HTML}{176C65}
\definecolor{boxbg}{HTML}{F2F6F8}
\hypersetup{colorlinks=true,linkcolor=ink,urlcolor=accent,citecolor=ink,
 pdftitle={Depth-Two Threshold Circuits for Andreev's Function: Sharper Bounds},
 pdfauthor={Christian Reitwiessner}}
\pagestyle{fancy}
\fancyhf{}
\fancyhead[L]{\small\color{ink}Threshold circuits for Andreev's function}
\fancyhead[R]{\small\thepage}
\renewcommand{\headrulewidth}{0.3pt}
\setlength{\parindent}{0pt}
\setlength{\parskip}{0.55em}
\setlength{\emergencystretch}{2em}
\setlist[itemize]{leftmargin=1.5em,itemsep=0.2em,topsep=0.2em}
\setlist[enumerate]{leftmargin=1.8em,itemsep=0.2em,topsep=0.2em}
\newtheorem{theorem}{Theorem}
\newtheorem{lemma}[theorem]{Lemma}
\newtheorem{corollary}[theorem]{Corollary}
\newtheorem{proposition}[theorem]{Proposition}
\newtheorem{fact}[theorem]{Fact}
\theoremstyle{remark}
\newtheorem*{remark}{Remark}
\newcommand{\bits}{\{0,1\}}
\newcommand{\ind}[1]{\bigl[\,#1\,\bigr]}
\newcommand{\sect}[1]{\section{\texorpdfstring{\color{ink}#1}{#1}}}
\newcommand{\subsect}[1]{\subsection*{\texorpdfstring{\color{ink}#1}{#1}}}
\newcommand{\smalltitle}[1]{\textbf{\color{ink}#1}\quad}
\newcommand{\THR}{\mathsf{THR}}
\newcommand{\cH}{\mathcal H}
\newtcolorbox{keybox}[1][]{breakable,colback=boxbg,colframe=ink!40,boxrule=0.4pt,
  arc=2pt,left=8pt,right=8pt,top=5pt,bottom=5pt,#1}
\begin{document}
\thispagestyle{empty}
{\color{accent}\small RESEARCH NOTE\par}
\vspace{0.4em}
{\raggedright\color{ink}\LARGE\bfseries Depth-Two Threshold Circuits for\\[0.15em]
Andreev's Function: Sharper Bounds\par}
\vspace{0.5em}
{\large Christian Reitwiessner\par}
{\small November 2026\par}
\vspace{0.8em}

\begin{keybox}
\textbf{Abstract.} Andreev's function on $n$ input bits reads a bit from a
table at an address given by the parities of blocks of input bits. For
depth-two threshold circuits we show:
\begin{itemize}
\item an exact circuit with $\tfrac34 n^2+1$ gates (previously $O(n^3/\log n)$
      by Kane and Williams~\cite{KW});
\item for agreement on a $\tfrac12+\varepsilon$ fraction of inputs, matching
      bounds of $\varepsilon^2n^{3/2}$ gates up to polylogarithmic factors,
      for $\varepsilon$ between about $\log n/\sqrt n$ and $\tfrac18$
      (previously only the lower bound $\varepsilon^3n^{3/2}/\log^3n$);
\item worst-case lower bounds of $\Omega(n^{3/2}/\log^2 n)$ gates and
      $\Omega(n^{5/2}/\log^{5/2}n)$ wires (previously $\log^3 n$ and
      $\log^{7/2}n$ in the denominators).
\end{itemize}
In depth three we improve the top fan-in bound for Andreev's function. For
the depth-three function $B_n$ of Kane and Williams we outline improved
lower bounds, saving one logarithm, together with an average-case bound that
was implicit in their work, for $B_n$ with explicit parameters. The lower
bounds for Andreev's function and the average-case circuits are formalized in
merged Lean code; the $\tfrac34n^2$ circuit, the depth-three results and the
results on $B_n$ are formalized in open, not yet reviewed pull requests; one
proposition has only a proof sketch checked numerically on small cases (see
Section~\ref{sec:lean}).
\end{keybox}

%===================================================================
\sect{Setting and results}

\smalltitle{Andreev's function.}
Fix integers $k,m\ge1$ and put $N=2^k$. The input consists of
\begin{itemize}
\item a \emph{table} $x$ of $N$ bits $x_c$, one for each address
      $c\in\bits^k$, and
\item $k$ \emph{address blocks}; block $i$ consists of the bits
      $a_{i,1},\dots,a_{i,m}$. Write $a$ for all address bits and
      $|a_i|$ for the number of ones in block $i$.
\end{itemize}
The \emph{address} $b\in\bits^k$ has coordinates $b_i=|a_i|\bmod2$, and
\[
A(x,a)=x_b .
\]
The input length is $n=N+km$. In the \emph{balanced case} $km=N$, so $n=2N$
and $k=\log_2(n/2)$. All asymptotic statements refer to the balanced case.

\smalltitle{Threshold circuits.}
For a statement $P$, the bracket $\ind{P}$ is $1$ if $P$ holds and $0$
otherwise. A \emph{threshold gate} with inputs $z_j\in\bits$, real weights
$\omega_j$ and threshold $\tau$ computes $\ind{\sum_j\omega_jz_j\ge\tau}$. In a \emph{depth-two threshold
circuit} ($\THR\circ\THR$) some \emph{bottom} gates read input bits, and a
single \emph{top} gate reads the bottom gates and input bits. Its
\emph{size} is its total number of gates, top gate included; its
\emph{wires} are its nonzero weights. A circuit
\emph{$\varepsilon$-agrees} with $A$ if it outputs $A(x,a)$ on at least a
$\tfrac12+\varepsilon$ fraction of all inputs.

\smalltitle{Results and the state of the art.}
Stronger lower bounds against $\THR\circ\THR$ are known for functions in the
much larger class $\mathsf E^{\mathsf{NP}}$; they are proved by the
algorithmic method~\cite{Tam,ACW,CTW}. For functions in $\mathsf P$, which
includes Andreev's function and $B_n$, the best known gate lower bounds are of
order $n^{3/2}/\mathrm{polylog}\,n$~\cite{KW,KKL}, as also noted
in~\cite{CTW}. Our bounds improve the explicit polylogarithmic factors and
the dependence on $\varepsilon$.
\begin{center}
\small
\renewcommand{\arraystretch}{1.25}
\setlength{\tabcolsep}{5pt}
\begin{tabular}{@{}lllll@{}}
\toprule
Function & Bound & Best previously known & This paper & \\
\midrule
\multicolumn{5}{@{}l}{\emph{Functions in $\mathsf P$}}\\
Andreev, exact & gates, upper & $O(n^3/\log n)$ \cite{KW} & $\tfrac34n^2+1$ & Thm.~\ref{thm:exact}\\
Andreev, exact & gates, lower & \begin{tabular}[t]{@{}l@{}}$\Omega(n^{3/2}/\log^3n)$ \cite{KW}\\ $n^{3/2}/(\log n)^{O(1)}$ \cite{KKL}$^\dagger$\end{tabular} & $\Omega(n^{3/2}/\log^2n)$ & Thm.~\ref{thm:lower}\\
Andreev, exact & wires, lower & $n^{5/2}/(\log n)^{O(1)}$ \cite{KKL}$^\dagger$ & $\Omega(n^{5/2}/\log^{5/2}n)$ & Thm.~\ref{thm:lower}\\
Andreev, $\varepsilon$-agree & gates, upper & none stated & $O(\varepsilon^2n^{3/2}\sqrt{\log n})$ & Thm.~\ref{thm:avgupper}\\
Andreev, $\varepsilon$-agree & gates, lower & $\Omega(\varepsilon^3n^{3/2}/\log^3n)$ \cite{KW} & $\Omega(\varepsilon^2n^{3/2}/\log^2n)$ & Thm.~\ref{thm:lower}\\
Andreev, $\varepsilon$-agree & wires, lower & $\Omega(\varepsilon^3n^{5/2}/\log^{7/2}n)$ \cite{KW} & $\Omega(\varepsilon^2n^{5/2}/\log^{5/2}n)$ & Thm.~\ref{thm:lower}\\
Andreev, depth 3 & top fan-in & $\tilde\Omega(n^{1/2}/s^{1/3})$, implicit in~\cite{KW} & $\min\{\tilde\Omega(n^{1/2}),\tilde\Omega(n^{3/4}/s^{1/2})\}$ & Cor.~\ref{cor:fanin}\\
$B_n$, depth 3, exact & gates, lower & $\Omega(n^{3/2}/\log^3n)$ \cite{KW} & $\Omega(n^{3/2}/\log^2n)$ & Thm.~\ref{thm:B}$^*$\\
$B_n$, depth 3, exact & wires, lower & $\Omega(n^{5/2}/\log^{7/2}n)$ \cite{KW} & $\Omega(n^{5/2}/\log^{5/2}n)$ & Thm.~\ref{thm:B}$^*$\\
$B_n$, $\varepsilon$-agree & gates, lower & $\tilde\Omega(\varepsilon n^{3/2})$, implicit in~\cite{KW} & $\Omega(\varepsilon n^{3/2}/\log^2n)$ & Thm.~\ref{thm:B}$^*$\\
\midrule
\multicolumn{5}{@{}l}{\emph{Functions in $\mathsf{E}^{\mathsf{NP}}$}}\\
some $F$, exact & gates, lower & \begin{tabular}[t]{@{}l@{}}$n^{2-\delta}$ \cite{Tam,ACW}\\ $n^{5/2-\delta}$ \cite{CTW}\end{tabular} & -- & \\
some $F$, $\varepsilon$-agree & wires, lower & $n^3/\log^{12}n$ \cite{Nag} (unreviewed) & -- & \\
\bottomrule
\end{tabular}
\end{center}
Notes on the table:
\begin{itemize}
\item All rows concern $\THR\circ\THR$ circuits unless they say depth~3.
      The first row is exact; the other rows are asymptotic in the
      balanced case.
\item The $\varepsilon$-agreement rows for Andreev's function hold for
      $3\sqrt2(\log_2n+1)/\sqrt n\le\varepsilon$, and the upper bound for
      $\varepsilon\le\tfrac18$. For smaller $\varepsilon$, a single
      majority gate over the table already agrees on a
      $\tfrac12+\Omega(1/\sqrt n)$ fraction.
\item In the top fan-in row, $s$ bounds the size of each depth-two
      subcircuit.
\item $^\dagger$For a variant of Andreev's function whose table is four times
      as long as the address part.
\item $^*$For $B_n$ with the parameters of Theorem~\ref{thm:B}; the
      $\varepsilon$-agreement row holds for $\varepsilon\ge2\cdot2^{-3p}$, about
      $n^{-3}$. Proved here as an outline with constants; the full proof is
      formalized (open pull request). Chen, Santhanam and
      Srinivasan~\cite[Thm.~9.3]{CSS} re-prove the Kane--Williams bound for
      $B_n$ with $n^{o(1)}$ losses; Kabanets, Kane and Lu~\cite[Thm.~1.2]{KKL}
      prove a related wire bound for a variant of $B_n$.
\item Because of the $\tfrac34n^2$ upper bound, no gate lower bound above
      $n^2$ holds for Andreev's function, so the $n^{5/2-\delta}$ bound
      of~\cite{CTW} cannot hold for it.
\end{itemize}
The remaining gap for Andreev's function is the worst case: between
$n^{3/2}/\log^2n$ and $\tfrac34 n^2$ gates.

\smalltitle{How to read this paper.}
Section~\ref{sec:tech} explains which techniques are new.
Sections~\ref{sec:upper}--\ref{sec:avgupper} (upper bounds) and
Section~\ref{sec:lower} (lower bounds) can be read independently;
Section~\ref{sec:depth3} uses Section~\ref{sec:lower}. Each lemma uses only
its own local variables plus the notation in the boxes, and each box says
where its notation is valid.

%===================================================================
\sect{Where the improvements come from}\label{sec:tech}

Each new result combines standard tools with one or two new observations.
This section is an overview; the terms are defined precisely later. Briefly:
a \emph{character} is a sign $(-1)^{S\cdot c}$ for a set $S$ of address
coordinates; a \emph{staircase} writes a function of an integer linear form as
a sum of threshold gates (Lemma~\ref{lem:ladder}); a \emph{restriction} fixes
all but one bit per address block, and a gate \emph{survives} if it still
depends on the remaining bits (Section~\ref{sec:lower}); for $B_n$, the
table is replaced by a \emph{seed} (Section~\ref{sec:depth3}).

\begin{center}
\small
\renewcommand{\arraystretch}{1.25}
\begin{tabular}{@{}P{0.39\textwidth}P{0.33\textwidth}P{0.2\textwidth}@{}}
\toprule
Technique & Status & Used for\\
\midrule
Fourier expansion of the table lookup & standard & Thms.~\ref{thm:exact},~\ref{thm:avgupper}\\
Telescoping a function of a linear form into threshold gates (``staircase'') & standard~\cite[Prop.~1]{HMPST}, \cite{HP} & Thms.~\ref{thm:exact},~\ref{thm:avgupper}\\
Cutting staircases to a window via concentration & \cite[Thm.~1.4]{KW}, for parity & Thm.~\ref{thm:avgupper}\\
\textbf{Grouping by character and removing the parity factor} (Lemmas~\ref{lem:lookup},~\ref{lem:parity}) & \textbf{new, to our knowledge}; similar in spirit to polarization identities & Thm.~\ref{thm:exact}\\
\textbf{Majority vote over a subcube} (sum over $S\subseteq\{1,\dots,k_0\}$) & \textbf{new in this context} & Thm.~\ref{thm:avgupper}\\
Random restrictions with Littlewood--Offord & \cite[Lemmas~1.1, 3.1]{KW} & Thm.~\ref{thm:lower}\\
Counting threshold functions over fixed features & \cite{RSO}, \cite[Thm.~2.3]{KW} & Thm.~\ref{thm:lower}\\
\textbf{Fixed-weight counting} (Lemma~\ref{lem:count}) & structural fact standard~\cite{KW,CSS}; \textbf{its use for counting new, to our knowledge} & one $\log$ in Thm.~\ref{thm:lower}\\
\textbf{Uniform-in-$q$ union bound with Cauchy--Schwarz} (Lemma~\ref{lem:avg}) & \textbf{new, to our knowledge}; \cite{KW,KKL,CSS} use Markov & $\varepsilon^2$ in Thm.~\ref{thm:lower}\\
Averaging over a family of circuits & \cite[Cor.~1.1]{KW} & Cor.~\ref{cor:fanin}\\
Nearly orthogonal vectors (Lemma~\ref{lem:gram}) & Johnson bound~\cite{Gur}; \cite[Lemma~6.1]{KW} & Thm.~\ref{thm:B}\\
\textbf{Fixing the circuit before the seed} & \textbf{new for $B_n$, to our knowledge}; a standard probabilistic move & Thm.~\ref{thm:B}\\
\bottomrule
\end{tabular}
\end{center}

\smalltitle{Upper bounds.}
Kane and Williams~\cite[Thm.~1.2]{KW} write $x_b$ as a threshold of about
$n^2/\log n$ parities and build each parity from its own depth-two
subcircuit of $O(n)$ gates. We group the parities that share the same
address character into one product of a linear form and a parity, then
remove the parity with an elementary identity (``removing the parity
factor'', Lemma~\ref{lem:parity}). This saves the factor $n/\log n$: a
generic simulation of a threshold of parities builds every parity separately,
whereas Lemma~\ref{lem:parity} lets all parities with the same address
character share one staircase. The average-case circuit keeps only the
characters of the first $k_0$ address bits, which turns the lookup into a
majority vote over a subcube.

\smalltitle{Lower bounds.}
We keep the restriction argument of Kane and Williams and change two counting
steps.
\begin{itemize}
\item \emph{Fixed-weight counting.} After the restriction, a surviving gate
      keeps its weights on the live bits; only its threshold changes. This
      fact is standard; it appears in the proof of~\cite[Lemma~1.1]{KW}, as
      equation~(9.1) of~\cite{CSS}, and in satisfiability algorithms for
      threshold circuits~\cite{CTW}. We use it for counting: a survivor costs
      $O(k+\log s)$ bits to describe, not the $\Theta(k^2)$ bits of an
      arbitrary threshold function, as in~\cite{KW,KKL,CSS}. This removes one
      logarithm. It works because the circuit is fixed before the random table
      is drawn; arguments that fix the hard input before the circuit, by
      counting~\cite{KW,CSS} or by incompressibility~\cite{KKL}, cannot do
      this.
\item \emph{Uniform union bound.} Kane and Williams bound the number of
      surviving gates with Markov's inequality, which costs a factor of
      $\varepsilon$. We bound the correlation for all survivor counts at
      once and average with Cauchy--Schwarz. This changes $\varepsilon^3$
      to $\varepsilon^2$.
\end{itemize}
For $B_n$, fixed-weight counting needs the circuit to be fixed before the
seed is chosen; Kane and Williams fix the seed first.

%===================================================================
\sect{An exact circuit with \texorpdfstring{$\tfrac34 n^2$}{3/4 n\^{}2} gates}\label{sec:upper}

\begin{keybox}
\textbf{Notation for Sections~\ref{sec:upper} and~\ref{sec:avgupper}.}
For $S\subseteq\{1,\dots,k\}$ and $c\in\bits^k$ write $S\cdot c=\sum_{i\in S}c_i$, and
\[
X_S=\sum_{c\in\bits^k}(-1)^{S\cdot c+x_c},\qquad
t_S=\sum_{i\in S}|a_i| .
\]
$X_S$ depends only on the table, is even, and lies in $[-N,N]$.
$t_S$ is the number of ones in the blocks in $S$, and lies in $[0,|S|m]$.
\end{keybox}

\smalltitle{Plan.}
The circuit computes the integer score $N\cdot(-1)^{x_b}$. Its top gate
tests whether the score is negative, which happens exactly when $x_b=1$.
Lemma~\ref{lem:lookup} writes the score as a sum of products,
Lemma~\ref{lem:parity} removes the products, and Lemma~\ref{lem:ladder}
turns the result into threshold gates.

\begin{lemma}[Lookup as a sum of products]\label{lem:lookup}
$\displaystyle\sum_{S\subseteq\{1,\dots,k\}}X_S\,(-1)^{t_S}=N\,(-1)^{x_b}.$
\end{lemma}
\begin{proof}
Since $|a_i|\equiv b_i\pmod2$, we have $(-1)^{t_S}=(-1)^{S\cdot b}$.
Exchanging the two sums, the left side is
\[
\sum_c(-1)^{x_c}\prod_{i=1}^k\bigl(1+(-1)^{c_i+b_i}\bigr),
\]
and the product is $N$ for $c=b$ and $0$ otherwise.
\end{proof}

Let $\varphi(\theta)=(-1)^\theta\,\theta$ for integers $\theta$.

\begin{lemma}[Removing the parity factor]\label{lem:parity}
If $X$ is an even integer and $t$ is an integer, then
$X(-1)^t=\varphi(X+t)-\varphi(t)$.
\end{lemma}
\begin{proof}
$X+t$ has the parity of $t$, so $\varphi(X+t)=(-1)^t(X+t)$.
\end{proof}

\begin{lemma}[Threshold staircase]\label{lem:ladder}
Let $L$ be an integer that always lies in $\{\ell,\dots,h\}$, and let $g$
be any function on the integers. Then
\[
g(L)=g(\ell)+\sum_{\theta=\ell+1}^{h}\bigl(g(\theta)-g(\theta-1)\bigr)\ind{L\ge\theta}.
\]
\end{lemma}
\begin{proof}
The brackets equal $1$ exactly for $\theta\le L$, and the sum telescopes.
\end{proof}

When $L$ is a linear form in the input bits, each $\ind{L\ge\theta}$ is one
threshold gate. Both $t_S$ and $X_S+t_S$ are such forms, since
$(-1)^{x_c}=1-2x_c$:
\begin{equation}\label{eq:form}
X_S+t_S=\sum_c(-1)^{S\cdot c}-2\sum_c(-1)^{S\cdot c}x_c+\sum_{i\in S}\sum_{j=1}^m a_{i,j}.
\end{equation}

\begin{theorem}\label{thm:exact}
For all $k,m\ge1$, Andreev's function has an exact depth-two threshold
circuit with $2N^2+kmN$ bottom gates and one top gate. In the balanced case
this is $\tfrac34n^2+1$ gates. Bottom weights lie in $\{0,\pm2\}$ on table
bits and in $\{0,1\}$ on address bits, and
top weights are odd integers of absolute value at most $2n-1$.
\end{theorem}
\begin{proof}
By Lemmas~\ref{lem:lookup} and~\ref{lem:parity},
\[
N(-1)^{x_b}=\sum_S\bigl(\varphi(X_S+t_S)-\varphi(t_S)\bigr).
\]
Apply Lemma~\ref{lem:ladder} to every term. $X_S+t_S$ ranges over
$[-N,\,N+|S|m]$, which needs $2N+|S|m$ gates, and $t_S$ over $[0,|S|m]$,
which needs $|S|m$ gates. Each index $i$ lies in half of the $N$ sets $S$,
so $\sum_S|S|=kN/2$, and the number of bottom gates is
$\sum_S(2N+2|S|m)=2N^2+kmN$. The top gate weights the gate
$\ind{L\ge\theta}$ by $\pm(\varphi(\theta)-\varphi(\theta-1))=\pm(-1)^\theta(2\theta-1)$.
This is odd, and since $-N<\theta\le N+km=n$, its absolute value is at most
$2n-1$.
\end{proof}
The term $S=\varnothing$ is $X_\varnothing$, a linear form in the table that the
top gate could read directly; this would save $2N$ gates. We keep the uniform
construction.
Since the bottom weights are small integers, the circuit is even a threshold
of majority gates ($\THR\circ\mathsf{MAJ}$, by repeating inputs). For
comparison, Kane and Williams also give a $\mathsf{MOD}_3\circ\mathsf{MOD}_2$
circuit of size $O(n^2)$~\cite[Thm.~1.2]{KW}, in a different basis.

\begin{keybox}[unbreakable,title={Generalization to other moduli (proof sketch, checked numerically on small cases, not formalized)},
  fonttitle=\bfseries\small,coltitle=ink,colbacktitle=boxbg]
\begin{proposition}\label{prop:modp}
Let $P$ be a prime, and change the address to $b_i=|a_i|\bmod P$, with a
table indexed by $c\in\{0,\dots,P-1\}^k$ (so $N=P^k$). Then the function has
an exact depth-two threshold circuit with at most $PN^2+(P-1)^2kmN$
bottom gates and real weights.
\end{proposition}
\emph{Proof sketch.} For $S,c\in\{0,\dots,P-1\}^k$ let $S\cdot c=\sum_iS_ic_i$,
$t_S=\sum_iS_i|a_i|$, and $Z_{S,\nu}=\#\{c: S\cdot c\equiv \nu \pmod P,\ x_c=1\}$.
Summing the characters of $\mathbb Z_P^k$ and taking real parts gives
\[
N x_b=\sum_{S}\sum_{\nu=0}^{P-1}Z_{S,\nu}\cos\bigl(2\pi(t_S-\nu)/P\bigr).
\]
The term $S=0$ is $\sum_cx_c$, which the top gate reads directly. For
$S\ne0$, the function $g_\nu(\theta)=\cos(2\pi(\theta-\nu)/P)$ has period $P$,
so $g_\nu(PZ+t)=g_\nu(t)$ for every integer $Z$. Hence, with
$\psi_\nu(\theta)=\theta\,g_\nu(\theta)$,
\[
Z_{S,\nu}\,g_\nu(t_S)=\tfrac1P\bigl(\psi_\nu(PZ_{S,\nu}+t_S)-\psi_\nu(t_S)\bigr),
\]
and $\sum_\nu\psi_\nu(t_S)=0$. For prime $P$ and $S\ne0$, exactly $N/P$
addresses $c$ have $S\cdot c\equiv\nu$, so $Z_{S,\nu}\le N/P$, and
$PZ_{S,\nu}+t_S$ ranges over an interval of length at most
$N+(P-1)m\cdot\#\{i:S_i\ne0\}$; its staircase needs that many gates. Summing
over $S\ne0$ and $\nu$, and using $\sum_S\#\{i:S_i\ne0\}=k(P-1)N/P$, gives the
bound. The top gate tests $Nx_b\ge N/2$. We checked the circuit gate by gate
for $P=3$ with $k\le3$ and $P=5$ with $k=1$ (and, outside the proposition,
for $P=4$ with the bound $Z_{S,\nu}\le N$); these are small cases only.
\end{keybox}

%===================================================================
\sect{Agreeing on most inputs with fewer gates}\label{sec:avgupper}

\begin{theorem}\label{thm:avgupper}
\leavevmode
\begin{enumerate}
\item For every $\delta\in(0,1)$, there is a depth-two threshold circuit with
      $O\bigl(N\sqrt{n\log(N/\delta)}\bigr)$ gates that agrees with Andreev's
      function on a $1-\delta$ fraction of inputs.
\item For every $\varepsilon$ with $1/\sqrt N\le\varepsilon\le\tfrac18$,
      there is one with $O\bigl(\varepsilon^2N\sqrt{n\log n}\bigr)$ gates that
      $\varepsilon$-agrees with Andreev's function.
\end{enumerate}
In the balanced case these are $O(n^{3/2}\sqrt{\log(n/\delta)})$ and
$O(\varepsilon^2n^{3/2}\sqrt{\log n})$ gates. Part~1 with
$\delta=\tfrac12-\varepsilon$ covers $\tfrac18<\varepsilon\le\tfrac12-n^{-O(1)}$.
\end{theorem}

\smalltitle{The circuit.}
Choose $k_0\in\{0,\dots,k\}$. The circuit is the one from
Theorem~\ref{thm:exact}, with two changes: it uses only the sets
$S\subseteq\{1,\dots,k_0\}$, and each staircase covers only a window of
values. The top gate still tests whether the score is negative. We need one
consequence of Lemma~\ref{lem:ladder}: if the sum in that lemma is taken only
over a window $\{\ell,\dots,h\}$, it still equals $g(L)$ whenever $L$ lies in
the window.

\begin{proof}
\emph{Fewer terms.} The computation in Lemma~\ref{lem:lookup}, restricted to
$S\subseteq\{1,\dots,k_0\}$, gives
\[
\sum_{S\subseteq\{1,\dots,k_0\}}X_S(-1)^{t_S}=2^{k_0}\sum_{c\in\mathcal Q}(-1)^{x_c},
\]
where $\mathcal Q$ is the set of the $V=2^{k-k_0}$ addresses that agree with
$b$ in the first $k_0$ coordinates. So without windows the circuit outputs a
majority vote of the $V$ table entries in $\mathcal Q$, one of which is
$x_b$.

\emph{The vote on a random table.} Let $V\ge2$, so $V$ is even, and let $Y$
be the sum of the signs $(-1)^{x_c}$ of the other $V-1$ entries; $Y$ is odd.
The circuit outputs $1$ exactly when $(-1)^{x_b}+Y<0$. If $x_b=1$ this is
correct exactly when $Y\le-1$; if $x_b=0$, exactly when $Y\ge-1$. Hence
\[
\Pr[\text{correct}]=\tfrac12\bigl(1+\Pr[Y=-1]\bigr)
=\tfrac12+\tfrac12\tbinom{V}{V/2}2^{-V}\ \ge\ \tfrac12+\frac1{2\sqrt{2V}},
\]
using $\binom{V-1}{V/2}=\tfrac12\binom{V}{V/2}$ and the standard bound
$\binom{V}{V/2}\ge2^V/\sqrt{2V}$. For $V=1$ the vote is exactly $x_b$.

\emph{Windows.} On a random input, $X_S$ is a sum of $N$ independent uniform
signs, and $t_S$ is a sum of $|S|m\le n$ independent uniform bits. Hoeffding's
inequality gives, for every $\lambda>0$,
\[
\Pr\bigl[|X_S|>\lambda\bigr]\le2e^{-\lambda^2/(2N)}\le2e^{-\lambda^2/(2n)},\qquad
\Pr\bigl[|t_S-|S|m/2|>\lambda\bigr]\le2e^{-2\lambda^2/n}\le2e^{-\lambda^2/(2n)}.
\]
Take $\lambda=\sqrt{2n\ln(4\cdot2^{k_0}/\delta)}$. A union bound over the two
events and the $2^{k_0}$ sets shows that, except with probability
$2^{k_0}\cdot2\cdot2e^{-\lambda^2/(2n)}=\delta$, every $t_S$ lies within
$\lambda$ of $|S|m/2$ and every $X_S+t_S$ within $2\lambda$ of $|S|m/2$. On
these windows, rounded outwards to integers, the staircases give exactly the
score of the unwindowed circuit. For each $S$ the two windows need at most
$4\lambda+2$ and $2\lambda+2$ gates, so the circuit has at most
$2^{k_0}(6\lambda+4)+1$ gates.

\emph{Part 1.} Take $k_0=k$, so $V=1$; the output is wrong only outside the
windows, which has probability at most $\delta$.

\emph{Part 2.} Let $V$ be the largest power of two with
$V\le1/(32\varepsilon^2)$. Since $\varepsilon\le\frac18$ we have $V\ge2$, and
since $\varepsilon\ge1/\sqrt N$ we have $V\le N$, so $k_0=k-\log_2V$ is
allowed. Because $V\le1/(32\varepsilon^2)$, the vote has advantage
$1/(2\sqrt{2V})\ge2\varepsilon$. Take $\delta=\varepsilon$. Then the agreement
is at least $\tfrac12+2\varepsilon-\varepsilon$. Because
$V>1/(64\varepsilon^2)$, the number of gates is
$O(2^{k_0}\lambda)=O\bigl((N/V)\sqrt{n\log(N/\varepsilon)}\bigr)
=O(\varepsilon^2N\sqrt{n\log n})$.
\end{proof}

The lower bounds of Section~\ref{sec:lower} show that Part~2, and Part~1 for
$\delta\ge n^{-O(1)}$, are tight up to polylogarithmic factors. The Lean
formalization of this theorem uses the older identity
$\varphi(X+t)+\varphi(X-t)=2X(-1)^t$, which only changes constant factors, and
carries the redundant hypothesis $km\le n$.

%===================================================================
\sect{Lower bounds}\label{sec:lower}

\begin{theorem}\label{thm:lower}
Let $C$ be a depth-two threshold circuit of size $s$ with $w$ wires that
$\varepsilon$-agrees with Andreev's function, where
$\varepsilon\ge3(k+2)/\sqrt N$. Then
\[
m\,N\,\varepsilon^2\le 36\,\bigl(k+2+\ln(s+1)\bigr)\cdot s\cdot\sqrt n
\quad\text{and}\quad
m^{3/2}N\varepsilon^2\le 36\,\bigl(k+2+\ln(w+1)\bigr)\cdot w .
\]
If moreover $m\le N$, the first factor can be replaced by $4(k+2)$, giving
$mN\varepsilon^2\le144(k+2)\,s\sqrt n$ and
$m^{3/2}N\varepsilon^2\le144(k+2)\,w$. In the balanced case,
$s=\Omega(\varepsilon^2n^{3/2}/\log^2n)$ and
$w=\Omega(\varepsilon^2n^{5/2}/\log^{5/2}n)$. Taking $\varepsilon=\tfrac12$,
allowed when $N\ge36(k+2)^2$, gives the worst-case bounds.
\end{theorem}

\smalltitle{Plan.}
We follow Kane and Williams~\cite{KW}. Fix all but one bit per block at
random. The restricted circuit computes a lookup in a \emph{random} table
indexed by the $k$ remaining bits, and only a few bottom gates still depend
on those bits (Lemma~\ref{lem:surv}). There are few such restricted circuits
(Lemma~\ref{lem:count}), and none of them correlates well with a random
table (Lemma~\ref{lem:avg}). Averaging over all restrictions proves the
theorem.

\begin{keybox}
\textbf{Notation for Sections~\ref{sec:lower} and~\ref{sec:depth3}.}
$C$ is a depth-two threshold circuit of size $s$ with $w$ wires.
A \emph{restriction} consists of a choice $\sigma$ of one \emph{live}
position in every block, and a \emph{fixing}: values for all other input
bits, table included. For a uniformly random restriction, every fixed bit
is uniform and independent. Write $u\in\bits^k$ for the live bits and
$o_i$ for the parity of the $m-1$ fixed bits of block $i$ (the
\emph{offset}). Then the address is $b=u\oplus o$, and the restricted
function is
\[
u\mapsto x_{u\oplus o}.
\]
A bottom gate \emph{survives} if it still depends on $u$. Let $U$ be the
number of surviving gates.
\end{keybox}

\begin{lemma}[Few survivors]\label{lem:surv}
For a uniformly random restriction, $\mathbb E[U]\le4s\sqrt n/m$. The
expected number of bottom gates that depend on at least two live bits is at
most $4w/m^{3/2}$. Both bounds hold for every fixed table, since they use
only the randomness of the fixed address bits.
\end{lemma}
These are the random restriction lemmas of Kane and
Williams~\cite[Lemmas~1.1 and~3.1]{KW}, with explicit constants.
Appendix~\ref{app:surv} proves both bounds; the proof includes small repairs
to the argument of~\cite{KW} (ties among weights, zero weights, and a strict
versus non-strict inequality). An alternative proof of the first lemma, up to
constants, is in~\cite[App.~A.3]{CTW}; \cite[Lemma~6.1]{KKL}
and~\cite[Lemma~9.1]{CSS} prove analogues with polylogarithmic and $n^{o(1)}$
losses.

\begin{fact}[Roychowdhury, Siu, Orlitsky~\cite{RSO}; see {\cite[Thm.~2.3]{KW}}]\label{fact:rso}
For any fixed functions $f_1,\dots,f_q$ on a finite set $\Omega$, at most
$(|\Omega|+1)^{q+1}$ functions on $\Omega$ have the form
$\ind{\sum_j\omega_jf_j\ge\tau}$.
\end{fact}

\begin{lemma}[Counting restricted circuits]\label{lem:count}
Fix $C$ and $\sigma$, and let $\cH_q$ be the set of functions
$u\mapsto C(\text{restricted})$ over all fixings with at most $q$
survivors. Then
\[
|\cH_q|\le\bigl((2^k+1)(s+1)\bigr)^q\,(2^k+1)^{q+k+1}.
\]
\end{lemma}
\begin{proof}
Non-surviving bottom gates are constants, and are absorbed into the top
threshold. A surviving gate is $\ind{\omega\cdot u\ge\tau}$, where $\omega$ is
its original weight vector restricted to the live positions; only $\tau$
depends on the fixing. For a fixed $\omega$, these functions are nested as
$\tau$ grows, so there are at most $2^k+1$ of them. Pad the list of survivors
with constant functions to length exactly $q$. Then each entry of the list is
one of at most $(2^k+1)(s+1)$ functions: choose one of the at most $s$ gates
or ``constant'', then the threshold. The top gate is a threshold function of
the $q$ entries and the $k$ live bits. By Fact~\ref{fact:rso} with
$|\Omega|=2^k$, there are at most $(2^k+1)^{q+k+1}$ such functions for each
list.
\end{proof}
Kane and Williams~\cite{KW}, and also~\cite{KKL,CSS}, count all threshold
functions of the live bits at this step, which costs about $k^2$ bits per
survivor instead of $O(k+\log s)$.

\begin{lemma}[Agreement bound for all survivor counts at once]\label{lem:avg}
Let $\alpha=\sqrt{N(k+2+\ln(s+1))}$ and $\beta=(k+2)\sqrt N$. Then $C$
agrees with Andreev's function on at most a
\[
\tfrac12+\frac{\alpha\sqrt{\mathbb E[U]}+\beta}{N}+2^{-(k+6)}
\]
fraction of inputs.
\end{lemma}
\begin{proof}
Fix $\sigma$ and the fixed address bits; the table remains uniformly random.
For one function $f$ of $u$ and $\Delta\ge0$, Hoeffding's inequality gives
\[
\Pr\bigl[f\text{ agrees with }u\mapsto x_{u\oplus o}\text{ on}\ge N/2+\Delta\text{ points}\bigr]
\le e^{-2\Delta^2/N}.
\]
Call the fixing \emph{bad} if some $q<N$ and some $f\in\cH_q$ violate this
with $\Delta_q=\alpha\sqrt q+\beta$. Using $\Delta_q^2\ge\alpha^2q+\beta^2$,
$\ln(2^k+1)\le k+1$ and Lemma~\ref{lem:count},
\begin{align*}
\Pr[\text{bad}]&\le\sum_{q\ge0}|\cH_q|\,e^{-2\Delta_q^2/N}
\le\sum_{q\ge0}e^{(2q+k+1)(k+1)+q\ln(s+1)-2q(k+2+\ln(s+1))-2(k+2)^2}\\
&=\sum_{q\ge0}e^{-2q-q\ln(s+1)+(k+1)^2-2(k+2)^2}
\le\sum_{q\ge0}e^{-2q-(k^2+6k+7)},
\end{align*}
which is at most $1.2\,e^{-(k^2+6k+7)}\le2^{-(k+6)}$.

On a good fixing with $U<N$ survivors, the agreement on the $N$ points $u$
is at most $N/2+\alpha\sqrt U+\beta$. This also holds for $U\ge N$, because
then $\alpha\sqrt U\ge\sqrt{2N}\sqrt N>N/2$. Averaging over all restrictions
and using $\mathbb E\sqrt U\le\sqrt{\mathbb E U}$ (Cauchy--Schwarz) proves
the lemma, since the average agreement over all restrictions equals the
agreement over all inputs.
\end{proof}
Kane and Williams~\cite{KW} (and similarly~\cite{KKL,CSS}) instead bound $U$
by Markov's inequality and then use a single survivor count. That loses a
factor of $\varepsilon$; the uniform bound with Cauchy--Schwarz does not.

\begin{proof}[Proof of Theorem~\ref{thm:lower}]
By hypothesis, $\beta/N=(k+2)/\sqrt N\le\varepsilon/3$ and
$2^{-(k+6)}\le\varepsilon/192$. So Lemma~\ref{lem:avg} gives
$\varepsilon/3\le\alpha\sqrt{\mathbb E[U]}/N$. Squaring and inserting
Lemma~\ref{lem:surv},
\[
\varepsilon^2N^2\le9\alpha^2\,\mathbb E[U]\le
9N\bigl(k+2+\ln(s+1)\bigr)\cdot\frac{4s\sqrt n}{m},
\]
which is the first bound.

For wires, count as survivors only the gates that depend on at least two
live bits. A gate depending on at most one live bit computes that bit, its
negation, or a constant, so it can be absorbed into the top gate, which reads
the live bits anyway. Every remaining survivor has at least two input wires,
so there are at most $w$ candidate gates, and Lemmas~\ref{lem:count}
and~\ref{lem:avg} hold with $s$ replaced by $w$. Now use the second bound of
Lemma~\ref{lem:surv}.

Finally, let $m\le N$. If $\ln(s+1)\le3(k+2)$, then $k+2+\ln(s+1)\le4(k+2)$.
Otherwise $s\ge e^{3(k+2)}-1>N^{5/2}\ge mN$, so the claim with $4(k+2)$ holds
trivially. The same argument works for $w$, since $N^{5/2}\ge m^{3/2}N$.
\end{proof}

\begin{remark}[Limits of the method; informal]
After a restriction with one live bit per block, the function
$u\mapsto x_{u\oplus o}$ is an OR of at most $N$ ANDs of literals, which is a
depth-two circuit with $N$ bottom gates. Gates such as $\ind{t_S\ge\theta}$
survive with probability $\Theta(k/\sqrt N)$. So no argument that counts
survivors of a single restriction can give more than about $n^{3/2}/\log n$,
one logarithm above Theorem~\ref{thm:lower}. Kane and Williams already observe
that their restriction lemma is tight for majority gates~\cite[Section~3]{KW}.
Rectangle methods (discrepancy, sign rank) cannot reach $n^{2-o(1)}$ either:
for every split of the inputs between two players, Andreev's function has a
protocol with $O(k)$ bits of communication. One player sends its partial
block parities, and the other then knows $b$ and either announces $x_b$ or
sends $b$ back. Sign rank is also known to be too weak against
$\THR\circ\THR$ in general~\cite{CM}.
\end{remark}

%===================================================================
\sect{Depth three}\label{sec:depth3}

\subsect{Top fan-in for Andreev's function}

A family of depth-two circuits $C_1,\dots,C_M$ is an
\emph{$\varepsilon$-approximate majority} for $A$ if, on every input, at
least a $\tfrac12+\varepsilon$ fraction of the $C_i$ output $A$.

\begin{corollary}\label{cor:fanin}
Let $C_1,\dots,C_M$ be depth-two threshold circuits of size at most $s$
each, and assume $m\le N$.
\begin{enumerate}
\item If they form an $\varepsilon$-approximate majority for $A$ with
      $\varepsilon\ge3(k+2)/\sqrt N$, then
      $mN\varepsilon^2\le144(k+2)\,s\sqrt n$.
\item If on every input strictly more than half of them output $A$, then
      either $M>\sqrt N/(6(k+2))$ or $mN\le576(k+2)\,M^2s\sqrt n$.
\end{enumerate}
In the balanced case, Part~2 says that a majority of $M$ depth-two circuits
of size $s$ computing $A$ needs
$M\ge\min\bigl\{\sqrt N/(6(k+2)),\ \tilde\Omega(n^{3/4}/\sqrt s)\bigr\}$.
The same holds for wires, with $\sqrt n$ omitted and $m$ replaced by
$m^{3/2}$.
\end{corollary}
\begin{proof}
Count the pairs (input, $i$) for which $C_i$ is correct. Every input
contributes at least $(\tfrac12+\varepsilon)M$ pairs, so some $C_i$ is
correct on at least a $\tfrac12+\varepsilon$ fraction of inputs. Apply
Theorem~\ref{thm:lower} (the version for $m\le N$) to that $C_i$. For
Part~2, strictly more than half correct means at least
$\tfrac12+\tfrac1{2M}$, so take $\varepsilon=1/(2M)$; the condition on $M$
is the hypothesis of Theorem~\ref{thm:lower}.
\end{proof}
Kane and Williams state the example ``$n^{1.1}$ gates require top fan-in
$n^{2/15}$'', which corresponds to $M\ge\tilde\Omega(n^{1/2}/s^{1/3})$
(implicit in~\cite[Cor.~1.1]{KW}). With Corollary~\ref{cor:fanin}, a majority
of depth-two circuits of size $n^{1.1}$ that computes $A$ needs top fan-in
$n^{1/5-o(1)}$. Part~2 assumes nothing about how the top majority is built;
it only uses that more than half of the subcircuits are correct.

\subsect{The function \texorpdfstring{$B_n$}{B\_n} (proof outline; formalized in Lean)}

Andreev's function has depth-three majority circuits of size
$O(n)$~\cite[Thm.~1.2]{KW}, so depth-three lower bounds need a harder
function. Kane and Williams~\cite[Section~6]{KW} use the following one.

\begin{keybox}
\textbf{Notation for this subsection.} $k$ and $m$ keep their roles as the
number of address blocks and the block length; $u$, $o$, $\sigma$ and $U$
are as in Section~\ref{sec:lower}, and the address is again $b=u\oplus o$.
The table is replaced by a \emph{seed} $y\in\bits^d$ and a list
$r_z\in\bits^d$, indexed by $z\in\bits^k$, that is \emph{$\eta$-biased}: for
every nonzero $e\in\bits^d$, $\bigl|\sum_z(-1)^{r_z\cdot e}\bigr|\le\eta2^k$.
The function is
\[
B(y,a)=r_b\cdot y\bmod2 ,
\]
and the input length is $n=d+km$. For $f:\bits^k\to\bits$ and
$v\in\{\pm1\}^{2^k}$, the \emph{correlation} is
$\mathrm{corr}(f,v)=2^{-k}\bigl|\sum_z(-1)^{f(z)}v_z\bigr|$; if $f$ agrees with
the sign pattern $v$ on a $\tfrac12+\delta'$ fraction of points, then
$\mathrm{corr}(f,v)\ge2\delta'$. In this subsection, $C$ and $s$ may refer to
the depth-three circuit and its size; this is stated where it applies.
\end{keybox}

Kane and Williams take $d=km$, $k=16\log_2d$, and an explicit list with
$\eta=O(d^{-7})$ from Alon, Goldreich, H{\aa}stad and Peralta~\cite{AGHP}. We
use the same shape with explicit parameters: $k=16p$ blocks, $d=km\le2^p$ and
$\eta\le2^{-7p}$. The powering construction of~\cite{AGHP} gives such lists,
with bias at most $(d-1)2^{-8p}\le2^{-7p}$, so the function is in $\mathsf P$.

\begin{lemma}[Nearly orthogonal vectors]\label{lem:gram}
Let $v_1,\dots,v_J\in\{\pm1\}^D$ satisfy $|\langle v_i,v_j\rangle|\le\eta D$
for $i\ne j$, and let $\xi\in\{\pm1\}^D$. If $|\langle\xi,v_i\rangle|\ge\gamma D$
for every $i$, then $J(\gamma^2-\eta)\le1$.
\end{lemma}
\begin{proof}
Let $v=\sum_i\pm v_i$, with signs chosen so that each
$\langle\xi,\pm v_i\rangle\ge\gamma D$. Then $\langle\xi,v\rangle\ge J\gamma D$
and $\|v\|^2\le JD+J^2\eta D$. Cauchy--Schwarz gives
$J^2\gamma^2D^2\le D(JD+J^2\eta D)$, which rearranges to the claim.
\end{proof}
This is a version of~\cite[Lemma~6.1]{KW}, and a form of the Johnson
bound~\cite[Ch.~3]{Gur}; the sharper form is $J\le(1-\eta)/(\gamma^2-\eta)$.
Apply it with $D=2^k$ and $v_y=\bigl((-1)^{r_z\cdot y}\bigr)_z$: for seeds
$y\ne y'$, $\langle v_y,v_{y'}\rangle$ is the sum of the character of $y+y'$
over the list, so $|\langle v_y,v_{y'}\rangle|\le\eta D$. Hence, if
$\gamma^2>\eta$, a fixed function of $z$ has correlation at least $\gamma$
with at most $1/(\gamma^2-\eta)$ of the vectors $v_y$.

\begin{fact}[Discriminator; {\cite{HMPST}}, {\cite[Prop.~9.5]{CSS}}]\label{fact:disc}
Let $h_1,\dots,h_M:\bits^k\to\bits$ and let $\kappa$ be an integer. If the
function $\ind{h_1+\dots+h_M\ge\kappa}$ agrees everywhere with a function whose
sign pattern $v$ satisfies $2^{-k}|\sum_zv_z|\le\eta$, then some $h_i$ has
$\mathrm{corr}(h_i,v)\ge(1-(M+1)\eta)/M$.
\end{fact}

\needspace{14\baselineskip}
\begin{theorem}\label{thm:B}
Let $p\ge1$, $k=16p$, $d=km\le2^p$, and let the list be $2^{-7p}$-biased.
Write $K_1=272p^2+54p+3$ and $K_2=272p^2+57p+3$.
\begin{enumerate}
\item Let $C$ consist of depth-two threshold circuits, possibly sharing bottom
      gates, with $s$ gates and $w$ wires (the top gate and its wires are not
      counted), and a top gate that outputs $1$ exactly when at least $\kappa$
      of $M$ selected outputs are $1$. The outputs may repeat and are counted
      with multiplicity, and $M\le d^3/4$. If $C$ computes $B$, then
      \[
      d\,(d-K_1)\le4864\,p^2\,s\sqrt{2d}\qquad\text{and}\qquad
      d^{3/2}(d-K_1)\le19456\,p^{5/2}\,w .
      \]
\item If $d>K_2$, every depth-two threshold circuit with $s$ gates and $w$
      wires agrees with $B$ on at most a
      \[
      \tfrac12+2^{-3p}+\frac{2432\,p^2\,s\sqrt{2d}}{d\,(d-K_2)}
      \quad\text{and at most a}\quad
      \tfrac12+2^{-3p}+\frac{9728\,p^{5/2}\,w}{d^{3/2}(d-K_2)}
      \]
      fraction of inputs.
\end{enumerate}
\end{theorem}
\emph{Reading the bounds.} Take $m=\lfloor2^p/(16p)\rfloor$; for $p\ge9$ this
gives $2^{p-1}<d\le2^p$, so $p=\lceil\log_2d\rceil$. Part~1 then gives
$\Omega(n^{3/2}/\log^2n)$ gates and $\Omega(n^{5/2}/\log^{5/2}n)$ wires, and
Part~2 gives agreement $\tfrac12+O(s\log^2n/n^{3/2})+O(n^{-3})$, that is,
$\varepsilon$-agreement needs $\Omega(\varepsilon n^{3/2}/\log^2n)$ gates when
$\varepsilon\ge2\cdot2^{-3p}$. The bounds are nontrivial only for $p\ge17$; at
$p=17$, where $d=130832$, Part~1 certifies only about $10$ gates. The
hypothesis $M\le d^3/4$ costs nothing for majority circuits in the sense of
Kane and Williams: if the subcircuits are distinct then $M\le s$, and when
$s>d^3/4$ Part~1 holds trivially. Chen, Santhanam and Srinivasan use the same
kind of fan-in restriction~\cite[Lemma~9.8]{CSS}.

\begin{proof}[Proof outline with constants]
Let $\gamma=2^{-3p}$. Kane and Williams~\cite[Cor.~6.2]{KW}, and
also~\cite{CSS,KKL}, fix one seed that is bad for all small circuits. We fix the
circuit first, which allows fixed-weight counting. We may assume $s<2^{4p}$
(respectively $w<2^{4p}$), since otherwise the bounds hold trivially; then
$\log_2(s+1)\le4p$. Note that $m=d/(16p)$.
\begin{enumerate}
\item \emph{Survivor budget.} Let $q_0$ be the largest integer $Q$ with
      $(36p+2)Q+K\le d$, where $K=K_1$ in Part~1 and $K=K_2$ in Part~2. Then
      $q_0+1>(d-K)/(36p+2)$.
\item \emph{Few bad seeds.} For each $\sigma$ and each offset $o$, take the
      restricted functions with at most $q_0$ survivors, shifted by $o$. Call a
      seed \emph{bad} if one of them has correlation at least $\gamma$ with
      $v_y$, and also call the zero seed bad. By Lemma~\ref{lem:count},
      $\log_2|\cH_q|\le q(36p+2)+(16p+1)^2$. There are $m^k\le2^{16p^2}$
      choices of $\sigma$ and $2^k$ offsets. By Lemma~\ref{lem:gram}, each
      function correlates with at most $1/(\gamma^2-\eta)\le2^{6p+1}$ seeds.
      Together,
      \[
      \Pr[\text{bad seed}]\le2^{(36p+2)q_0+272p^2+54p+2-d}.
      \]
      With $K=K_1$ this is at most $\tfrac12$; with $K=K_2$ it is at most
      $2^{-3p-1}=\gamma/2$.
\item \emph{Few survivors.} Lemma~\ref{lem:surv} gives
      $\mathbb E[U]\le4s\sqrt{2d}/m=64p\,s\sqrt{2d}/d$, the expectation being
      over the restriction and the seed together. By Markov's inequality and
      step~1,
      \[
      \Pr[U>q_0]\le\frac{64p(36p+2)\,s\sqrt{2d}}{d(d-K)}
      \le\frac{2432\,p^2\,s\sqrt{2d}}{d(d-K)} .
      \]
      For wires, $\mathbb E[U]\le4w/m^{3/2}=256p^{3/2}w/d^{3/2}$, and
      $256p^{3/2}(36p+2)\le9728p^{5/2}$.
\item \emph{Part 2.} On a good seed with at most $q_0$ survivors, the circuit
      agrees with $B$ on at most a $\tfrac12(1+\gamma)$ fraction of the
      restricted inputs. So the agreement is at most
      $\tfrac12(1+\gamma)+\Pr[\text{bad seed}]+\Pr[U>q_0]\le\tfrac12+\gamma+\Pr[U>q_0]$.
\item \emph{Part 1.} Suppose $\Pr[U>q_0]<\tfrac12$. Together with step~2, some
      restriction has a good seed and at most $q_0$ survivors. On it, $C$
      computes exactly the function of $u$ whose sign pattern is
      $u\mapsto(v_y)_{u\oplus o}$. A good seed is nonzero, so this target has
      bias at most $\eta$. By
      Fact~\ref{fact:disc}, some subcircuit has correlation at least
      $(1-(M+1)\eta)/M>\gamma$ with it, since $M\le d^3/4\le2^{3p}/4$. That
      contradicts the choice of the seed. Hence $\Pr[U>q_0]\ge\tfrac12$, which
      by step~3 gives the factor $4864=2\cdot2432$, and $19456=2\cdot9728$ for
      wires.\qedhere
\end{enumerate}
\end{proof}
The formal proof (Section~\ref{sec:lean}) follows these steps.

Part~2 is not essentially new: Kane and Williams' proof
of~\cite[Cor.~6.2]{KW} combined with Markov's inequality gives the same form
with an extra logarithm, although they do not state it. They do say that they
cannot prove correlation bounds of the form $n^{-\omega(1)}$ for $B_n$, and
that their approach cannot go below correlation $1/\sqrt n$ for circuits with
about $n$ gates~\cite[Section~6]{KW}. Part~2 agrees: for $s\approx n$ it
gives about $\log^2n/\sqrt n$. What is new is the saved logarithm, obtained
by fixing the circuit first.

%===================================================================
\sect{Formalization status}\label{sec:lean}

The Lean~4 proofs are in the Complexitylib library,
\url{https://github.com/cslib-community/complexitylib}. Merged results are at
commit \texttt{242879e}; the others are in open pull requests that have not
yet been independently reviewed. The formal proofs contain no
\texttt{sorry}, and the library's continuous integration checks that they use
only the standard axioms.
\begin{center}
\small
\renewcommand{\arraystretch}{1.2}
\begin{tabular}{@{}P{0.29\textwidth}P{0.25\textwidth}P{0.39\textwidth}@{}}
\toprule
Result & Status & Where\\
\midrule
Theorem~\ref{thm:exact} & formalized (open PR) & PR \#2813, \texttt{Andreev.lean}\\
Proposition~\ref{prop:modp} & proof sketch; numerical check on small cases & --\\
Theorem~\ref{thm:avgupper} & formalized (merged) & \texttt{AndreevWindowed.lean}, \texttt{AndreevAverage.lean}\\
Lemma~\ref{lem:surv} & formalized (merged) & \texttt{RandomRestriction.lean}, \texttt{ThresholdCounting.lean}\\
Lemmas~\ref{lem:count},~\ref{lem:avg}, Thm.~\ref{thm:lower} & formalized (merged) & \texttt{ThresholdCounting.lean}, \texttt{AndreevLowerBound.lean}\\
Corollary~\ref{cor:fanin} & formalized (open PR) & PR \#2814, \texttt{AndreevApproxMajority.lean}\\
Lemma~\ref{lem:gram}, Theorem~\ref{thm:B} & formalized (open PR) & PR \#2816, \texttt{KaneWilliamsB.lean}, \texttt{NearlyOrthogonal.lean}\\
\bottomrule
\end{tabular}
\end{center}
The formal statements are explicit finite inequalities; the asymptotic
corollaries are derived in the text. Where constants are given, the formal
versions have the same or slightly better constants. Some details:
\begin{itemize}
\item Theorem~\ref{thm:exact}: the formal statement bounds the bottom
      weights by $2$ in absolute value; the finer description of the weights
      is not formalized.
\item Theorem~\ref{thm:avgupper}: the formal version is the windowed circuit
      with free window parameters, built with the older identity
      $\varphi(X+t)+\varphi(X-t)=2X(-1)^t$; the choices made in the proof of
      Part~2 are on paper.
\item Corollary~\ref{cor:fanin}: the top majority is modelled as a hypothesis
      on the number of correct subcircuits.
\item Theorem~\ref{thm:B}: in the Lean code the seed length $d$ is called $n$
      and the block length $m$ is called $\ell$. The circuit is a depth-two
      circuit with $M$ outputs (whose size and wire count are $s$ and $w$),
      followed by the gate ``at least $\kappa$ of the $M$ outputs''. The formal
      Part~2 uses $K_2$ and Part~1 uses $K_1$, as stated.
\end{itemize}

\smalltitle{Acknowledgements.}
The results, the formalization and this note were developed with AI
assistance.

\appendix
\renewcommand{\thesection}{\Alph{section}}
\sect{Proof of Lemma~\ref{lem:surv}}\label{app:surv}

\begin{fact}[Erd\H{o}s~\cite{Erdos}]\label{fact:lo}
Let $z_1,\dots,z_R$ be independent uniform bits, let $\omega_1,\dots,\omega_R$
satisfy $|\omega_j|\ge\mu>0$, and let $I$ be a half-open interval of length
$\mu$. Then
$\Pr\bigl[\sum_j\omega_jz_j\in I\bigr]\le\binom{R}{\lfloor R/2\rfloor}2^{-R}\le1/\sqrt R$.
\end{fact}

Fix one bottom gate $\ind{\sum_j\omega_jz_j\ge\tau}$, where $j$ runs over
all input positions. Let $\Lambda$ be the set of address positions with
$\omega_j\ne0$, so $|\Lambda|$ is at most the number of the gate's input wires.
For $j\in \Lambda$ let $\rho(j)$ be the number of $j'\in \Lambda$ with
$|\omega_{j'}|\ge|\omega_j|$, including $j$ itself. Fix the table bits; they
only shift the threshold.

\begin{lemma}\label{lem:onerestr}
Fix the live positions. Over the random fixed address bits, the gate is
non-constant on the live bits with probability at most
$2\sum_{j}\rho(j)^{-1/2}$, the sum over the live positions $j\in \Lambda$.
\end{lemma}
\begin{proof}
Let $\Phi$ be the contribution of the fixed bits. The restricted gate is
$\ind{\sum_{j\text{ live}}\omega_ju_j\ge\tau-\Phi}$, and it is non-constant
exactly when $\tau-\Phi$ lies in a half-open interval of length
$\sum_{j\text{ live}}|\omega_j|$. Split this interval into consecutive
half-open pieces of lengths $|\omega_j|$, one for each live $j\in \Lambda$.

For a live $j\in \Lambda$, let $R_j$ be the number of \emph{fixed} positions
$j'\in \Lambda$ with $|\omega_{j'}|\ge|\omega_j|$. If $R_j\ge\rho(j)/4$ for every live
$j\in \Lambda$, then Fact~\ref{fact:lo}, applied to these $R_j$ positions and
conditioned on all other fixed bits, bounds the probability of piece $j$ by
$R_j^{-1/2}\le2\rho(j)^{-1/2}$, and summing over the pieces proves the lemma.

Otherwise some live $j$ has $R_j<\rho(j)/4$, so more than $3\rho(j)/4$ live
positions $j'\in \Lambda$ have $|\omega_{j'}|\ge|\omega_j|$, and each of them has
$\rho(j')\le\rho(j)$. Then the claimed bound is at least
$\tfrac34\rho(j)\cdot2\rho(j)^{-1/2}\ge\tfrac32$, which exceeds any
probability.
\end{proof}

\begin{proof}[Proof of Lemma~\ref{lem:surv}]
Each position of a block is live with probability $1/m$. Sorting the values
$\rho(j)$, the $t$-th smallest is at least $t$. So the gate survives with
probability at most
\[
\frac2m\sum_{j\in \Lambda}\rho(j)^{-1/2}\le\frac2m\sum_{t=1}^{|\Lambda|}t^{-1/2}
\le\frac{4\sqrt{|\Lambda|}}{m}\le\frac{4\sqrt n}{m}.
\]
Summing over the at most $s$ bottom gates gives the first bound.

A gate that depends on at least two live bits is non-constant and has at
least two live positions in $\Lambda$. Positions in different blocks are live
independently, so for $j\in \Lambda$,
\[
\Pr[\,j\text{ and another position of }\Lambda\text{ are live}\,]
\le\min\Bigl\{\frac1m,\frac{|\Lambda|}{m^2}\Bigr\}\le\frac{\sqrt{|\Lambda|}}{m^{3/2}}.
\]
By Lemma~\ref{lem:onerestr}, the gate depends on at least two live bits with
probability at most
\[
2\sum_{j\in \Lambda}\frac{\sqrt{|\Lambda|}}{m^{3/2}}\,\rho(j)^{-1/2}
\le\frac{2\sqrt{|\Lambda|}}{m^{3/2}}\cdot2\sqrt{|\Lambda|}=\frac{4|\Lambda|}{m^{3/2}}.
\]
Summing over the bottom gates, and using that $\sum|\Lambda|$ is at most the number
of wires $w$, gives the second bound.
\end{proof}

The formal proofs are the Lean theorems
\texttt{card\_not\_forcedToConst\_ltfEval\_le} (Lemma~\ref{lem:onerestr}, there
with all non-live positions random),
\texttt{sum\_circuitUnforcedConstCount\_blockLiveSet\_le} and
\texttt{sum\_circuitUnforcedAtMostOneCount\_blockLiveSet\_le}.

\renewcommand{\refname}{\texorpdfstring{\color{ink}References}{References}}
\begin{thebibliography}{99}\setlength{\itemsep}{0pt}\footnotesize

\bibitem{AGHP} N.~Alon, O.~Goldreich, J.~H{\aa}stad, and R.~Peralta.
\emph{Simple constructions of almost $k$-wise independent random variables}.
Random Structures \& Algorithms 3(3):289--304, 1992.
\bibitem{ACW} J.~Alman, T.~M. Chan, and R.~Williams.
\emph{Polynomial Representations of Threshold Functions and Algorithmic
Applications}. FOCS 2016, pages 467--476.
\bibitem{CM} A.~Chattopadhyay and N.~S. Mande.
\emph{Weights at the Bottom Matter When the Top is Heavy}.
\href{https://eccc.weizmann.ac.il/report/2017/083/}{ECCC TR17-083}, 2017.
\bibitem{CSS} R.~Chen, R.~Santhanam, and S.~Srinivasan.
\emph{Average-Case Lower Bounds and Satisfiability Algorithms for Small
Threshold Circuits}. Theory of Computing 14(9):1--55, 2018.
\bibitem{CTW} L.~Chen, A.~Tal, and Y.~Wang.
\emph{Super-quadratic Lower Bounds for Depth-2 Linear Threshold Circuits}.
STOC 2026. \href{https://eccc.weizmann.ac.il/report/2026/039/}{ECCC TR26-039}.
\bibitem{Erdos} P.~Erd\H{o}s. \emph{On a lemma of Littlewood and Offord}.
Bull.\ Amer.\ Math.\ Soc.\ 51:898--902, 1945.
\bibitem{Gur} V.~Guruswami. \emph{List Decoding of Error-Correcting Codes}.
LNCS 3282, Springer, 2004.
\bibitem{HP} K.~A. Hansen and V.~V. Podolskii. \emph{Polynomial threshold
functions and Boolean threshold circuits}. Information and Computation
240:56--73, 2015.
\bibitem{HMPST} A.~Hajnal, W.~Maass, P.~Pudl\'ak, M.~Szegedy, and
G.~Tur\'an. \emph{Threshold circuits of bounded depth}. J.\ Comput.\ Syst.\
Sci.\ 46(2):129--154, 1993.
\bibitem{KKL} V.~Kabanets, D.~M. Kane, and Z.~Lu.
\emph{A Polynomial Restriction Lemma with Applications}. STOC 2017.
\href{https://eccc.weizmann.ac.il/report/2017/026/}{ECCC TR17-026}.
\bibitem{KW} D.~M. Kane and R.~Williams.
\emph{Super-Linear Gate and Super-Quadratic Wire Lower Bounds for Depth-Two
and Depth-Three Threshold Circuits}. STOC 2016.
\href{https://arxiv.org/abs/1511.07860}{arXiv:1511.07860v1}.
\bibitem{Nag} D.~Nag. \emph{Almost-Everywhere Near-Cubic Wire Lower Bounds for
$\mathsf{SYM}\circ\THR$ and $\THR\circ\THR$}. Preprint,
\href{https://eccc.weizmann.ac.il/report/2026/167/}{ECCC TR26-167}, 2026.
\bibitem{RSO} V.~P. Roychowdhury, K.-Y. Siu, and A.~Orlitsky.
\emph{Neural models and spectral methods}. In Theoretical Advances in Neural
Computation and Learning, pages 3--36, 1994.
\bibitem{Tam} S.~Tamaki. \emph{A Satisfiability Algorithm for Depth Two
Circuits with a Sub-Quadratic Number of Symmetric and Threshold Gates}.
\href{https://eccc.weizmann.ac.il/report/2016/100/}{ECCC TR16-100}, 2016.
\end{thebibliography}
\end{document}
